% Irodov Part V §5.1 original 46--57. Source printed pp.194--196,
% original diagrams Fig.5.10--5.11; source answers printed 326--327.
\opticsproblem{46}{194}
开普勒望远镜原角放大率 $\Gamma=15$。
望远镜内外充满水，但筒长不变，为重新实现望远镜无穷远调整，只更换物镜。
目镜玻璃折射率为 $n=1.50$。求新角放大率。
\exsolution
原目镜在空气中焦距为 $f_e$，
原物镜 $f_o=\Gamma f_e$，
筒长 $\ell=(\Gamma+1)f_e$。
浸入折射率 $n_w\simeq1.33$ 的水后，
同一玻璃薄目镜焦距
$f_{e,w}=f_e\,n_w(n-1)/(n-n_w)$，
记其倍率为
$K=n_w(n-1)/(n-n_w)$。
新物镜调整为
$f_{o,w}=\ell-f_{e,w}
=(\Gamma+1-K)f_e$。
故
\[
\boxed{\Gamma_w=
\frac{f_{o,w}}{f_{e,w}}
=(\Gamma+1)\frac{n-n_w}{n_w(n-1)}-1
\simeq3.1}.
\]
\excheck 系统筒长保持不变意味着\emph{更换物镜}，不是简单将两焦距按同一个比例缩放。

\opticsproblem{47}{194}
望远镜物镜口径
$D=6.0$cm，
人眼瞳孔直径
$d_0=3.0$mm。
忽略光损耗，问角放大率 $\Gamma$ 最多可以多大，
以使视网膜像照度不低于裸眼观察？
\exsolution
由望远镜光束的拉格朗日不变量，
物镜通光口径 $D$ 在目镜后形成的出瞳直径
$d_{\rm exit}=D/\Gamma$。
若出瞳比眼瞳小，则有部分裸眼可接受光锥面积被压缩，
会降低扩展目标的视网膜照度。
不降低照度的条件
$d_{\rm exit}\ge d_0$，所以
\[
\boxed{\Gamma\le\frac D{d_0}
=\frac{60\,\mathrm{mm}}{3.0\,\mathrm{mm}}
=20}.
\]
\excheck 该约束适用于扩展目标、忽略损耗与无杂散光；不能不加限制套用于单点星等探测。

\opticsproblem{48}{194}
显微镜物镜和目镜的光焦度分别是
$100$D、$20$D，
当前角放大率 $\Gamma=50$。
两镜间距增加 $2.0$cm，问新放大率。
\exsolution
物镜焦距 $f_o=1/100$ m $=1.0$cm，
目镜焦距 $f_e=1/20$m $=5.0$cm。
显微镜近轴放大率
$\Gamma\simeq\Delta\,\ell_0/(f_of_e)$，
其中 $\Delta$ 为光学管长，
$\ell_0=25$cm 为标准明视距离。
间距增加 $\delta\Delta=2.0$cm，
增量
\[
\delta\Gamma\simeq\frac{\ell_0\delta\Delta}
{f_of_e}
=\frac{25(2.0)}{(1.0)(5.0)}
=10.
\]
故
\[
\boxed{\Gamma_{\rm new}\simeq60}.
\]
\excheck 这里保持物镜、目镜焦距不变，假设中间实像依然位于目镜可调整的成像条件下。

\opticsproblem{49}{194}
显微镜物镜数值孔径
$\sin\alpha=0.12$，
眼瞳直径
$d_0=4.0$mm。
确定放大率
(a) 使进入眼瞳的光束直径恰等于眼瞳直径；
(b) 使视网膜像照度不依赖显微镜放大率（受物镜框限光时）。
\exsolution
小孔径近轴模型中，经过显微镜后到达眼瞳的光束直径约
$d_{\rm exit}\simeq2\alpha\ell_0/\Gamma$，
$\ell_0=25$cm 是明视距离。
(a) 令
$d_{\rm exit}=d_0$，
有
\[
\boxed{\Gamma_*=\frac{2\alpha\ell_0}{d_0}
\simeq\frac{2(0.12)(0.25)}{0.004}
=15}.
\]
(b) 当 $\Gamma\le\Gamma_*$，
出射光束直径不小于眼瞳，瞳孔限光使视网膜像亮度近似不再受放大率影响，
\[
\boxed{\Gamma\le15}.
\]
\excheck 若 $\Gamma>15$，显微镜出瞳小于眼瞳，改变放大率会改变视网膜接收的照度与角面积关系。

\opticsproblem{50}{194}
一只薄对称双凸玻璃透镜两面曲率半径模
$R=7.50$cm，
一侧为空气、一侧为水。
求两个主平面、物方及像方焦点、节点位置。
取玻璃折射率 $n\simeq1.50$、水折射率
$n_w\simeq1.33$。
\exsolution
薄透镜的两个主平面重合于其中心面 H。
空气入射、水出射时总光焦度
\[
\Phi=\frac{n-1}R+\frac{n_w-n}{-R}
=\frac{2n-n_w-1}{R}
\simeq\frac{0.67}{0.075}
\simeq8.9\,\mathrm D.
\]
物方、像方焦距分别为
\[
\boxed{f=-\frac1\Phi\simeq-11\,\mathrm{cm}},
\quad
\boxed{f'=\frac{n_w}{\Phi}\simeq+15\,\mathrm{cm}}.
\]
入射、出射介质折射率不同，两个节点不再与主点重合；
二节点重合于由光学节点法则确定的共轴位置，
其偏离透镜平面的有向量为
\[
\boxed{x_N=f'+f\simeq3.7\,\mathrm{cm}}
\]
（在水侧）。
\excheck 薄透镜两\emph{主平面}重合不意味着两个\emph{节点}也恰在透镜中心；需考虑两侧折射率差。

\opticsproblem{51}{195}
原书图5.10给出三个共轴系统，
按图示作出焦点和主平面：
(a) 由焦距 $f_1=1.5a$、
$f_2=-1.5a$ 的会聚与发散薄透镜组成的远摄系统，间距为 $a$；
(b) 由两只会聚薄透镜
$f_1=1.5a$、
$f_2=0.5a$ 组成，间距仍为 $a$；
(c) 厚凹凸透镜两面光焦度为
$\Phi_1=+50$D、
$\Phi_2=-50$D，
两面间距 $d=4$cm、内部折射率 $n=1.5$。
\exsolution
对相距 $d$、中间介质折射率 $n$ 的两折射面，
总光焦度
$\Phi=\Phi_1+\Phi_2-(d/n)\Phi_1\Phi_2$；
主平面相对前后面的有向距离
\[
\boxed{x=\frac{d\Phi_2}{n\Phi}},\qquad
\boxed{x'=-\frac{d\Phi_1}{n\Phi}}.
\]
(a) 空气中
$\Phi_1=2/(3a)$、
$\Phi_2=-2/(3a)$，
$\Phi=4/(9a)$，
所以
$\boxed{f=9a/4}$，
$\boxed{x=-3a/2,\ x'=-3a/2}$。
(b)
$\Phi_1=2/(3a)$、
$\Phi_2=2/a$，
$\Phi=4/(3a)$，
于是
$\boxed{f=3a/4}$，
$\boxed{x=3a/2,\ x'=-a/2}$。
(c)
$\Phi=-(d/n)(+50)(-50)=66.7$D，
$\boxed{f=1/\Phi=1.5\,\mathrm{cm}}$，
$\boxed{x=-2.0\,\mathrm{cm},\ x'=-2.0\,\mathrm{cm}}$。
依这些有向位置分别在原图轴线上竖起主平面，
并从 H、H' 向物像两侧量取焦距即可。
\excheck 主平面可落在实体镜组\emph{外部}；其有向位置不能仅凭镜组外形猜测。

\opticsproblem{52}{195}
按图5.11(a)(b)(c)所给光轴、焦点、主平面、部分已知共轭点或一条入射—出射光线，用几何作图求缺失的
(a) 像方焦点 F' 与主平面 H'；
(b) 点 S 的像点 S'；
(c) 前后焦点 F、F' 与后主平面 H'。
\exsolution
(a) 由已知平行轴入射光线确定主平面上等高对应点，
其折射后的延长线与轴线交于像方焦点 F'；
在给定任一非轴光线的入、出射延长线上截取相同高度点，分别作垂直轴线的前后主平面。
(b) 由 S 发出一条平行轴的辅助光线，在前主平面 H 上标记高度；
在后主平面 H' 相同高度处连接 F' 得其出射线；
再取由 S 指向前焦点 F 的第二辅助光线，其经系统后平行主轴出射；
两条出射线（必要时用反向延长线）的交点为 S'。
(c) 已知图中一对入射—出射光线的延长线，
在光轴方向上等高对应构造主平面；
平行轴入射光在后焦平面处汇聚定义 F'，
反向追迹平行轴出射光确定 F，
从而补齐图示标注。
\excheck 对图解题，解法的可检验对象是两条主光线的几何共轭条件；任何具体位置都须依据原图给定相对坐标作图，本稿不伪造未给出的长度。

\opticsproblem{53}{195}
按原书列出的四种共轴顺序
(a) F—S—H—H'—F'；
(b) H—S—F'—F—H'；
(c) H'—S—F'—F—H；
(d) F'—H'—S—H—F，
分别作出物点 S 的像点 S'。
\exsolution
统一利用物像共轭的两主光线：
第一条从 S 平行光轴到前主平面 H，
在后主平面 H' 的同一高度处转为指向 F' 的出射线；
第二条从 S 沿指向前焦点 F 的方向入射，
在后主平面 H' 同一高度点转为平行光轴的出射线。
两出射线（或其反向延长线）的交点定义像点 S'。
任一情况下都可用有向牛顿共轭式代数核对：
\[
\boxed{xx'=ff'},
\]
其中 $x$ 是 S 到物方焦点 F 的有向距离，
$x'$ 为像点 S' 到像方焦点 F' 的有向距离。
(a) S 位于 F 与 H 间，出射线反向延长线相交，得到物方虚像；
(b)—(d) 按相应给定点序，从 F 量取 $x$，利用 $x'=ff'/x$ 在 F' 的正确有向一侧标出 S'，再以两主光线交点复核。
\excheck 这四组符号顺序涉及非通常次序的主平面和焦点，必须使用有向距离；不能简单地把所有光路解释成普通薄会聚透镜。

\opticsproblem{54}{195}
远摄镜组由前会聚透镜
$\Phi_1=+10$D 与后发散透镜
$\Phi_2=-10$D 组成。
求：(a) 间距 $d=4.0$cm 时的总焦距、主平面位置；
(b) 使总焦距 $f$ 与前透镜至后焦点距离
$\ell$ 的比值最大时的间距及最大比值。
\exsolution
镜组位于空气中，
总光焦度
$\Phi=\Phi_1+\Phi_2-d\Phi_1\Phi_2=100d$
（$d$ 用米），故
$\boxed{f=1/(100d)}$。
(a) $d=0.04$m 时
$\Phi=4$D，
$\boxed{f=25\,\mathrm{cm}}$。
主平面位置（相对对应透镜顶点）：
\[
\boxed{x=\frac{d\Phi_2}\Phi=-10\,\mathrm{cm}},
\qquad
\boxed{x'=-\frac{d\Phi_1}\Phi=-10\,\mathrm{cm}}.
\]
(b) 后主平面距后透镜有向距离 $-0.10$m；
后焦点在从前透镜向右 $d-0.10+f$ 处，
所以
$\ell=d-0.10+1/(100d)$。
比值
\[
\frac f\ell=\frac1{100d^2-10d+1}.
\]
分母在 $d=0.05$m 最小，故
\[
\boxed{d_{\max}=5.0\,\mathrm{cm}},
\quad
\boxed{(f/\ell)_{\max}=\frac43}.
\]
\excheck 远摄镜组的有效焦距可以大于从前镜到焦点的机械长度；这正是远摄压缩结构的目的。

\opticsproblem{55}{195--196}
厚弯月形玻璃透镜的凸面曲率半径
$R_1=10.0$cm、凹面
$R_2=5.0$cm，厚
$d=3.0$cm。
求前后主平面和前、后焦点的位置。取玻璃 $n=1.50$、外为空气，并规定光从凸面入射。
\exsolution
前面有向曲率半径 $R_1=+0.10$m，
后面凹面按光从左向右传播的曲率中心同在右侧，
$R_2=+0.05$m。
两面光焦度
$\Phi_1=(n-1)/R_1=+5$D、
$\Phi_2=(1-n)/R_2=-10$D。
厚透镜总焦度
\[
\Phi=\Phi_1+\Phi_2-\frac d n\Phi_1\Phi_2
=5-10-\frac{0.03}{1.5}(5)(-10)
=\boxed{-4.0\,\mathrm D}.
\]
因此有效像方焦距
$\boxed{f'=-25\,\mathrm{cm}}$，
物方焦距 $f=+25$cm。
两主平面相对于凸面顶点和凹面顶点的有向位置为
\[
\boxed{x=\frac{d\Phi_2}{n\Phi}=+5.0\,\mathrm{cm}},
\quad
\boxed{x'=-\frac{d\Phi_1}{n\Phi}=+2.5\,\mathrm{cm}}.
\]
按上述有向约定先标主平面，随后在各主平面向物方、像方量取有向焦距即可确定焦点。
\excheck 弯月透镜总体为发散型；
主平面并不必处于厚透镜实体内部，前后面的法线和曲率号必须保持一致。

\opticsproblem{56}{196}
两只共轴薄透镜焦距为
$f_1,f_2$，相距 $d$。
希望用单只薄透镜代替整个镜组，
使任意物距处的横向放大率与原系统相同。
求替代透镜焦距与其相对第一只镜的位置。
\exsolution
两薄透镜在空气中的光焦度
$\Phi=1/f_1+1/f_2-d/(f_1f_2)$。
等效焦距
\[
\boxed{f_{\rm eq}=\frac{f_1f_2}{f_1+f_2-d}}.
\]
要在任意物距保持横向放大率，不能只匹配光焦度；
还须把替代透镜的主平面置于原镜组的前主平面 H。
该主平面距第一只镜的有向距离
\[
\boxed{x=\frac{d(1/f_2)}{\Phi}
=\frac{f_1d}{f_1+f_2-d}}.
\]
\excheck 若 $f_1+f_2-d=0$，镜组为无焦系统，
替代焦距趋于无穷，此时无法用普通有限焦距单薄透镜描述。

\opticsproblem{57}{196}
一只薄对称双凸玻璃透镜表面曲率半径模
$R=38$cm；
垂直光轴的平面镜在透镜后方
$\ell=12$cm 处，
透镜与镜间充满水。
求光从空气侧入射、经水区镜面反射再返回的等效光焦度。
\exsolution
取玻璃折射率 $n=1.50$、水折射率
$n_w\simeq1.33$。
从空气入射、经透镜进入水中时单次透镜光焦度
\[
\Phi'=\frac{n-1}{R}+
\frac{n_w-n}{-R}
=\frac{2n-n_w-1}{R}.
\]
光往返两次经过透镜，并穿越水层至平面镜返回，
其等效反射光焦度是
\[
\boxed{\Phi_{\rm eff}
=2\Phi'-\frac{2\ell}{n_w}(\Phi')^2}.
\]
代 $R=0.38$m、$\ell=0.12$m，
$\Phi'=(3-1.33-1)/0.38\simeq1.76$D，
得到
\[
\boxed{\Phi_{\rm eff}\simeq
2(1.76)-\frac{0.24}{1.33}(1.76)^2
\simeq3.0\,\mathrm D}.
\]
\excheck 平面镜本身不提供曲率光焦度，
但光往返水层的传播距离进入等效光焦度的二次修正项。
